%% parts/p03/finite-model-theory.tex --- P3 finite-model-theory boundary.
%% Includes the finite-model comparison and the computable-categoricity entry.
%% The EF method itself is in qe-categoricity-games-supplement.tex.

\chapter[Finite Model Theory]{Finite Model Theory: Asymptotics and Complexity}
\label{sec:p03-finite-model-theory}

\noindent\textit{Reader note.}  This is a separately authored finite-model
survey following the exact inherited chapters.  Its finite-structure,
zero--one-law, and descriptive-complexity boundaries follow
\cite{EbbinghausFlum1995,Libkin2004}; they do not replace inherited first-order
semantics or game arguments.

This chapter fixes a finite purely relational signature \(\tau\), all of whose
relation symbols have positive arity, with no constant or function symbols,
unless a result says otherwise.  A labeled
\(n\)-element \(\tau\)-structure has universe
\([n]=\{0,\ldots,n-1\}\); ``uniformly random'' means uniformly distributed
over all such labeled structures.  These restrictions are part of the
zero--one law, not expository conveniences.

\section{The zero--one law}
\label{sec:p03-zero-one-law}

\begin{hbtheorem}{thm:p03-fol-zero-one}{cited=Fagin1976}{FO zero--one law}
For every first-order \(\tau\)-sentence \(\varphi\), if \(\mathcal A_n\) is
uniformly random among the labeled \(\tau\)-structures with universe \([n]\),
then
\[
  \lim_{n\to\infty}\Pr[\mathcal A_n\entails\varphi]
  \quad\text{exists and belongs to }\{0,1\}.
\]
\end{hbtheorem}

This is the theorem of Glebski\u\i--Kogan--Liogon'ki\u\i--Talanov and,
independently, Fagin \cite{GlebskiiEtAl1969,Fagin1976,Libkin2004}.  It is an
asymptotic statement: it neither decides truth in each finite cardinality nor
decides finite validity.  Common enrichments by a built-in order or successor
can destroy the law.  Constants and function symbols require a separately
specified probability space and lie outside this formulation.

For this paragraph only, switch from the chapter's uniform distribution on
all structures to \(G(n,1/2)\): the uniform labeled irreflexive symmetric
graphs on \([n]\), equivalently the model in which each unordered pair is an
edge independently with probability \(1/2\).  In this probability space, the
sentences whose probability tends to one form the complete theory of the
countable random graph from
Theorem~\ref{thm:p03-random-graph}.  For each fixed pair of finite disjoint
vertex sets of sizes \(u,v\), the probability that no remaining vertex
realizes the prescribed adjacency pattern is
\[
 \bigl(1-2^{-(u+v)}\bigr)^{n-u-v}.
\]
There are at most \(n^{u+v}\) ordered choices of the two sets, so the union
bound makes the probability of \emph{any} failure at those fixed sizes at
most
\(n^{u+v}(1-2^{-(u+v)})^{n-u-v}\), which tends to zero.  Hence every fixed
extension axiom holds with probability tending to one.  Those axioms are complete by the same
back-and-forth argument that proves uniqueness of the countable random graph.
Thus the random graph is the countable model of the limiting theory, not the
limit of the finite graphs as one set-theoretic structure
\cite{GlebskiiEtAl1969,Fagin1976,Rado1964}.

\section{Existential second-order logic and NP}
\label{sec:p03-fagin-theorem}

\begin{hbdefinition}{def:p03-existential-second-order}{Existential second-order sentence}
An \emph{existential second-order sentence} over \(\tau\) has the form
\[
  \exists R_1\cdots\exists R_m\,\theta,
\]
where the \(R_i\) are new relation variables of fixed finite arities and
\(\theta\) is first-order in \(\tau\cup\{R_1,\ldots,R_m\}\).  This syntactic
class is introduced because it can guess finite relations serving as an NP
certificate; it is strictly different in standing from the existential
\emph{first-order} fragment used in preservation theorems.
\index{existential second-order logic}
\end{hbdefinition}

\begin{hbtheorem}{thm:p03-fagin}{cited=Fagin1974}{Fagin's theorem}
Over finite relational structures, an isomorphism-closed property belongs to
\(\mathsf{NP}\) if and only if it is definable by an existential second-order
sentence.
\end{hbtheorem}

The forward direction encodes an accepting polynomial-time computation by
existentially quantified relations; the reverse direction guesses those
relations and checks the first-order matrix in polynomial time on the fixed
finite vocabulary.  The complete theorem fixes the representation of finite
structures and is cited from \cite{Fagin1974,Libkin2004}; it does not identify
NP with existential first-order logic.

\section{Finite versus unrestricted consequence}
\label{sec:p03-finite-unrestricted}

The single public Trakhtenbrot statement is the exact inherited
Theorem~\ref{thm:trakhtenbrot}: for every finite relational signature
containing a binary relation symbol, finite satisfiability is c.e. undecidable
and finite validity is not c.e.  This chapter uses that theorem without
creating a second statement identity \cite{Trakhtenbrot1950}.

\begin{hbexample}{def:p03-finite-compactness-failure}{Failure of compactness over finite structures}
In a signature containing \(<\), let \(\Gamma\) consist of the strict-linear-
order axioms together with, for every \(m\geq1\), a sentence asserting at least
\(m\) distinct elements.  Every finite subset of \(\Gamma\) has a finite
model, but \(\Gamma\) itself has no finite model.  Ordinary first-order
compactness still gives a model---necessarily infinite---so the failure is in
the restriction to finite structures.
\end{hbexample}

\begin{landscape}
\begingroup\small\setlength{\tabcolsep}{4pt}
\begin{longtable}{@{}p{1.22in}p{1.62in}p{2.25in}p{2.06in}@{}}
\caption{Finite-model-theory results and unrestricted contrasts.}
\label{tab:p03-finite-model-theory}\\
\toprule
Result & Exact scope & Lookup conclusion & Boundary with unrestricted semantics \\
\midrule\endfirsthead
\toprule
Result & Exact scope & Lookup conclusion & Boundary with unrestricted semantics \\
\midrule\endhead
First-order zero--one law & Fixed finite positive-arity relational vocabulary; uniform labeled structures on \([n]\) & Every first-order sentence has limiting probability \(0\) or \(1\) \cite{GlebskiiEtAl1969,Fagin1976}. & Common built-in order/successor enrichments can destroy the law; it does not decide finite validity. \\
\addlinespace
Fagin's theorem & Isomorphism-closed properties of finite relational structures & \(\mathsf{NP}\) is exactly existential second-order definability \cite{Fagin1974}. & The quantifiers over certificate relations are second-order, not existential first-order quantifiers. \\
\addlinespace
Trakhtenbrot's theorem & Finite relational signature containing a binary relation & Finite satisfiability is c.e. undecidable; finite validity is not c.e. \cite{Trakhtenbrot1950}. & Unrestricted first-order validity is c.e. by soundness and completeness of an effective calculus. \\
\addlinespace
Finite compactness & Consequence restricted to finite models & Compactness fails, as the unbounded-finite-order example shows. & Ordinary compactness permits the infinite model that finite semantics forbids. \\
\bottomrule
\end{longtable}
\endgroup
\end{landscape}

Together, the table prevents three invalid transfers: ordinary compactness does
not become finite compactness; effective axiomatizability of unrestricted
validity does not transfer to finite validity; and the zero--one law does not
supply a finite-validity algorithm.  The EF criterion and its parity and
connectivity applications appear in the preceding chapter and provide the
principal inexpressibility method for this boundary.

\section{A computable-model-theory entry point}
\label{sec:p03-computable-model-theory}

\begin{hbdefinition}{def:p03-computable-presentation}{Computable presentations}
For a countable structure \(A\), a presentation with domain \(\mathbb N\) is
\emph{computable} when its atomic diagram is decidable.  A computable copy is
an isomorphic structure with such a presentation.  This name is needed because
the existence of an abstract countable copy says nothing about whether its
atomic facts can be decided.  It is a property of a concrete presentation,
not finite-model decidability or a property of every isomorphic copy.  The
distinction is used in the computable-categoricity definition below.
\index{computable presentation}
\index{computable model theory}
\end{hbdefinition}

\begin{hbdefinition}{def:p03-computably-categorical}{Computably categorical structure}
A computable structure \(A\) is \emph{computably categorical} when every
computable copy \(B\cong A\) admits a computable isomorphism \(A\to B\).
The adjective is needed because ordinary categoricity supplies an isomorphism
but says nothing about its computational complexity.  It is a property of one
computable isomorphism type relative to computable copies, not a cardinal
spectrum of a theory.  It differs from having a computable presentation, which
requires only one copy, and is used in the effective back-and-forth theorem.
\index{computable categoricity}
\end{hbdefinition}

\begin{hbtheorem}{thm:p03-q-computably-categorical}{proved=annex:p3.finite.1}{The rational order is computably categorical}
Every two computable presentations of the countable dense linear order without
endpoints are computably isomorphic.  In particular,
\((\mathbb Q,<)\) is computably categorical.
\end{hbtheorem}

For a broader treatment of computable structures, see
\textcite{AshKnight2000}.
